% !TEX program = tectonic
% =============================================================================
%  AULA 10 — Drones, Satélites e Redes Não Terrestres (NTN)
%  Curso: Redes sem Fio  |  Profa. Anelise Munaretto
%  FANET, UAV, LEO satélite, 3GPP NTN.
%  Compilar: tectonic aula10.tex
% =============================================================================
\documentclass[aspectratio=169,11pt]{beamer}
% =============================================================================
%  PREÁMBULO COMÚN para as aulas de Redes sem Fio (Beamer + TikZ)
%  Incluido com % =============================================================================
%  PREÁMBULO COMÚN para as aulas de Redes sem Fio (Beamer + TikZ)
%  Incluido com % =============================================================================
%  PREÁMBULO COMÚN para as aulas de Redes sem Fio (Beamer + TikZ)
%  Incluido com \input{preamble.tex} en cada aula.
%  Paleta de cores e estilos são definidos aqui uma única vez.
% =============================================================================

\usepackage[T1]{fontenc}
\usepackage{amsmath,amssymb,mathtools}
\usepackage{graphicx}
\usepackage{tikz}
\usetikzlibrary{positioning,arrows.meta,calc,backgrounds,decorations.pathmorphing,fit,shadows,shapes.symbols,shapes.geometric}
\usepackage{pgfplots}
\pgfplotsset{compat=1.16}
\usepackage{tabularx,booktabs,multirow,array}
\usepackage[normalem]{ulem}
\usepackage{hyperref}

% ---------- Paleta de cores própria ----------
\definecolor{aneliseAzul}{RGB}{20,60,110}
\definecolor{aneliseVerde}{RGB}{30,140,70}
\definecolor{aneliseLaranja}{RGB}{215,120,20}
\definecolor{aneliseGris}{RGB}{90,90,90}

% ---------- Tema ----------
\usetheme{Madrid}
\usecolortheme{whale}
\setbeamercolor{structure}{fg=aneliseAzul}
\setbeamercolor{title}{fg=aneliseAzul}
\setbeamercolor{frametitle}{fg=aneliseAzul}
\setbeamercolor{block title}{fg=aneliseAzul,bg=aneliseGris!12!white}
\setbeamercolor{block body}{bg=aneliseGris!7!white}
\hypersetup{colorlinks=true,linkcolor=aneliseAzul,urlcolor=aneliseVerde}

% ---------- Comandos auxiliares ----------
\newcommand{\kurose}{Kurose \& Ross, \emph{Computer Networking: a Top-Down Approach}}
\newcommand{\forouzan}{Forouzan, \emph{Data Communications and Networking}}
\newcommand{\stallings}{Stallings, \emph{Wireless Communications \& Networks}}
\newcommand{\tanenbaum}{Tanenbaum, \emph{Computer Networks}}
\newcommand{\rfc}[1]{\href{https://www.rfc-editor.org/rfc/rfc#1.txt}{RFC~#1}}
% -----------------------------------------------------------------------------
%  Estilos TikZ comuns para desenhar redes (posicionamento relativo)
% -----------------------------------------------------------------------------
\tikzset{
  host/.style={draw,rounded corners=2mm,fill=aneliseAzul!15,inner sep=1.5mm,font=\scriptsize},
  rt/.style={draw,rounded corners=1mm,fill=aneliseLaranja!25,inner sep=1mm,font=\scriptsize},
  nube/.style={draw,cloud,aspect=2.2,fill=aneliseGris!15,inner sep=1.5mm,font=\scriptsize},
  el/.style={draw,rounded corners=1.2mm,fill=aneliseGris!18,inner sep=0.8mm,font=\scriptsize},
  enl/.style={thick},
  enlA/.style={thick,-Stealth},
} en cada aula.
%  Paleta de cores e estilos são definidos aqui uma única vez.
% =============================================================================

\usepackage[T1]{fontenc}
\usepackage{amsmath,amssymb,mathtools}
\usepackage{graphicx}
\usepackage{tikz}
\usetikzlibrary{positioning,arrows.meta,calc,backgrounds,decorations.pathmorphing,fit,shadows,shapes.symbols,shapes.geometric}
\usepackage{pgfplots}
\pgfplotsset{compat=1.16}
\usepackage{tabularx,booktabs,multirow,array}
\usepackage[normalem]{ulem}
\usepackage{hyperref}

% ---------- Paleta de cores própria ----------
\definecolor{aneliseAzul}{RGB}{20,60,110}
\definecolor{aneliseVerde}{RGB}{30,140,70}
\definecolor{aneliseLaranja}{RGB}{215,120,20}
\definecolor{aneliseGris}{RGB}{90,90,90}

% ---------- Tema ----------
\usetheme{Madrid}
\usecolortheme{whale}
\setbeamercolor{structure}{fg=aneliseAzul}
\setbeamercolor{title}{fg=aneliseAzul}
\setbeamercolor{frametitle}{fg=aneliseAzul}
\setbeamercolor{block title}{fg=aneliseAzul,bg=aneliseGris!12!white}
\setbeamercolor{block body}{bg=aneliseGris!7!white}
\hypersetup{colorlinks=true,linkcolor=aneliseAzul,urlcolor=aneliseVerde}

% ---------- Comandos auxiliares ----------
\newcommand{\kurose}{Kurose \& Ross, \emph{Computer Networking: a Top-Down Approach}}
\newcommand{\forouzan}{Forouzan, \emph{Data Communications and Networking}}
\newcommand{\stallings}{Stallings, \emph{Wireless Communications \& Networks}}
\newcommand{\tanenbaum}{Tanenbaum, \emph{Computer Networks}}
\newcommand{\rfc}[1]{\href{https://www.rfc-editor.org/rfc/rfc#1.txt}{RFC~#1}}
% -----------------------------------------------------------------------------
%  Estilos TikZ comuns para desenhar redes (posicionamento relativo)
% -----------------------------------------------------------------------------
\tikzset{
  host/.style={draw,rounded corners=2mm,fill=aneliseAzul!15,inner sep=1.5mm,font=\scriptsize},
  rt/.style={draw,rounded corners=1mm,fill=aneliseLaranja!25,inner sep=1mm,font=\scriptsize},
  nube/.style={draw,cloud,aspect=2.2,fill=aneliseGris!15,inner sep=1.5mm,font=\scriptsize},
  el/.style={draw,rounded corners=1.2mm,fill=aneliseGris!18,inner sep=0.8mm,font=\scriptsize},
  enl/.style={thick},
  enlA/.style={thick,-Stealth},
} en cada aula.
%  Paleta de cores e estilos são definidos aqui uma única vez.
% =============================================================================

\usepackage[T1]{fontenc}
\usepackage{amsmath,amssymb,mathtools}
\usepackage{graphicx}
\usepackage{tikz}
\usetikzlibrary{positioning,arrows.meta,calc,backgrounds,decorations.pathmorphing,fit,shadows,shapes.symbols,shapes.geometric}
\usepackage{pgfplots}
\pgfplotsset{compat=1.16}
\usepackage{tabularx,booktabs,multirow,array}
\usepackage[normalem]{ulem}
\usepackage{hyperref}

% ---------- Paleta de cores própria ----------
\definecolor{aneliseAzul}{RGB}{20,60,110}
\definecolor{aneliseVerde}{RGB}{30,140,70}
\definecolor{aneliseLaranja}{RGB}{215,120,20}
\definecolor{aneliseGris}{RGB}{90,90,90}

% ---------- Tema ----------
\usetheme{Madrid}
\usecolortheme{whale}
\setbeamercolor{structure}{fg=aneliseAzul}
\setbeamercolor{title}{fg=aneliseAzul}
\setbeamercolor{frametitle}{fg=aneliseAzul}
\setbeamercolor{block title}{fg=aneliseAzul,bg=aneliseGris!12!white}
\setbeamercolor{block body}{bg=aneliseGris!7!white}
\hypersetup{colorlinks=true,linkcolor=aneliseAzul,urlcolor=aneliseVerde}

% ---------- Comandos auxiliares ----------
\newcommand{\kurose}{Kurose \& Ross, \emph{Computer Networking: a Top-Down Approach}}
\newcommand{\forouzan}{Forouzan, \emph{Data Communications and Networking}}
\newcommand{\stallings}{Stallings, \emph{Wireless Communications \& Networks}}
\newcommand{\tanenbaum}{Tanenbaum, \emph{Computer Networks}}
\newcommand{\rfc}[1]{\href{https://www.rfc-editor.org/rfc/rfc#1.txt}{RFC~#1}}
% -----------------------------------------------------------------------------
%  Estilos TikZ comuns para desenhar redes (posicionamento relativo)
% -----------------------------------------------------------------------------
\tikzset{
  host/.style={draw,rounded corners=2mm,fill=aneliseAzul!15,inner sep=1.5mm,font=\scriptsize},
  rt/.style={draw,rounded corners=1mm,fill=aneliseLaranja!25,inner sep=1mm,font=\scriptsize},
  nube/.style={draw,cloud,aspect=2.2,fill=aneliseGris!15,inner sep=1.5mm,font=\scriptsize},
  el/.style={draw,rounded corners=1.2mm,fill=aneliseGris!18,inner sep=0.8mm,font=\scriptsize},
  enl/.style={thick},
  enlA/.style={thick,-Stealth},
}

\title[Redes sem Fio]{Redes sem Fio\\ Drones, Satélites e\\ Redes Não Terrestres}
\subtitle{Aula 10}
\author[Anelise Munaretto]{Profa. Anelise Munaretto}
\institute[Curso]{Redes sem Fio --- Ciência da Computação}
\date{2026}

\begin{document}

\begin{frame}[plain]
  \titlepage
\end{frame}

\begin{frame}{Objetivos de aprendizagem}
  Ao final desta aula, você será capaz de:
  \begin{itemize}
    \item \textbf{explicar} o papel das constelações LEO e \textbf{diferenciar} FANET de MANET;
    \item \textbf{identificar} as categorias de NTN (LEO, MEO, GEO, HAPS) e comparar suas latências;
    \item \textbf{descrever} os desafios de Doppler, \emph{handover} e energia em drones e satélites;
    \item \textbf{calcular} o período orbital e o tempo de visibilidade de um satélite LEO;
    \item \textbf{analisar} a integração das NTN ao 5G (3GPP Rel.~17/18).
  \end{itemize}
  \begin{alertblock}{Pré-requisitos}
    Aulas 3 (redes ad hoc) e 7 (5G e redes móveis).
  \end{alertblock}
\end{frame}

\section{FANET}

\begin{frame}{FANET --- Flying Ad-Hoc Network}
  \begin{columns}
    \begin{column}{0.55\textwidth}
      \textbf{Rede formada por VANTs (UAVs)}
      \begin{itemize}
        \item Comunicação multiplataforma (drone a drone)
        \item Nós com alta mobilidade (até 100~km/h)
        \item Topologia 3D em mudança constante
        \item Desafios:
          \begin{itemize}
            \item link budget variável
            \item handover frequente
            \item energia limitada (bateria)
            \item roteamento 3D
          \end{itemize}
      \end{itemize}
    \end{column}
    \begin{column}{0.45\textwidth}
      \textbf{Aplicações}
      \begin{itemize}
        \item vigilância e monitoramento
        \item agricultura de precisão
        \item entrega de pacotes
        \item busca e salvamento
        \item mapeamento 3D
      \end{itemize}
    \end{column}
  \end{columns}
  \begin{center}
    \begin{tikzpicture}[font=\scriptsize,
        d/.style={draw,circle,minimum size=5mm,fill=aneliseAzul!15,inner sep=0.4mm},
      ]
      \node[d](d1){};
      \node[d,right=1.6cm of d1](d2){};
      \node[d,above=1.2cm of d2](d3){};
      \node[d,right=1.6cm of d2](d4){};
      \node[d,below=1.2cm of d4](d5){};
      \draw[thick,aneliseLaranja] (d1) -- (d2) -- (d3);
      \draw[thick,aneliseLaranja] (d2) -- (d4) -- (d5);
      \draw[dashed,thick,aneliseGris] (d3) to[bend left=20] (d4);
    \end{tikzpicture}
    \\[1mm] {\scriptsize \textbf{Figura 1:} topologia FANET em 3D com enlaces drone a drone e topologia em mudança constante.}
  \end{center}
  \note{FANET é um caso particular de MANET com mobilidade muito maior e topologia 3D; reforçar que a alta velocidade (até 100 km/h) e o link budget variável tornam o handover e o roteamento mais desafiadores que em redes veiculares.}
\end{frame}

\section{NTN}

\begin{frame}{Redes Não Terrestres (NTN) --- 3GPP}
  \begin{block}{Definição}
    Redes que utilizam plataformas aerotransportadas ou espaciais (satélites LEO/MEO/GEO, HAPS, UAVs).
  \end{block}
  \begin{columns}
    \begin{column}{0.5\textwidth}
      \textbf{Categorias}
      \begin{itemize}
        \item \textbf{LEO}: 300--1500~km, latência 20--40~ms
        \item \textbf{MEO}: 7000--25000~km, latência 50--80~ms
        \item \textbf{GEO}: 35~786~km, latência $>$ 250~ms
        \item \textbf{HAPS}: 18--25~km (estratosfera)
      \end{itemize}
    \end{column}
    \begin{column}{0.5\textwidth}
      \textbf{Integração 5G}
      \begin{itemize}
        \item 3GPP Rel.\ 17 (2022): suporte básico NTN
        \item Rel.\ 18 (2024): NTN NR, IoT-NTN
        \item Desafios: Doppler elevado, atraso de propagação, handover entre satélites
      \end{itemize}
    \end{column}
  \end{columns}
\end{frame}

\begin{frame}{Constelações LEO: Starlink, OneWeb, Kuiper}
  \begin{itemize}
    \item \textbf{Starlink} (SpaceX): $>$ 7000 satélites, banda Ku/Ka/E, laser ISL, latência 20--40~ms
    \item \textbf{OneWeb}: ~650 satélites, banda Ku, parceria com Eutelsat
    \item \textbf{Project Kuiper} (Amazon): ~3200 satélites, banda Ka
    \item \textbf{Telesat Lightspeed}: LEO, foco em backhaul corporativo
  \end{itemize}
  \begin{center}
    \begin{tikzpicture}[font=\scriptsize,
        t/.style={draw,rectangle,fill=aneliseLaranja!20,rounded corners=1mm,inner sep=1mm},
        s/.style={draw,circle,fill=aneliseAzul!15,inner sep=0.5mm},
      ]
      \node[t](terra){Terminal};
      \node[s,above=1.6cm of terra](s1){Sat};
      \node[s,above=1.2cm of s1](s2){Sat};
      \node[s,above=1.2cm of s2](s3){Sat};
      \node[t,above=1.6cm of s3](gw){Gateway};
      \draw[thick,aneliseVerde] (terra) -- (s1) -- (s2) -- (s3) -- (gw);
      \node[right=1mm of s2,font=\tiny]{enlace laser ISL};
    \end{tikzpicture}
    \\[1mm] {\scriptsize \textbf{Figura 2:} enlace satelital LEO com ISL (Inter-Satellite Link) laser entre terminais terrestres.}
  \end{center}
  \note{Starlink já ultrapassou 7000 satélites e usa ISL a laser para reduzir a dependência de gateways terrestres; comparar latência LEO (20--40 ms) com GEO (mais de 250 ms).}
\end{frame}

\section{Desafios}

\begin{frame}{Desafios de comunicação com drones e NTN}
  \begin{columns}
    \begin{column}{0.5\textwidth}
      \textbf{Mobilidade 3D}
      \begin{itemize}
        \item roteamento geográfico adaptativo (GPSR, GRP)
        \item predição de trajetória com ML
        \item handover suave entre drones/satélites
      \end{itemize}
      \textbf{Energia}
      \begin{itemize}
        \item otimização de rota para economizar bateria
        \item energy harvesting (painéis solares em UAV/HAPS)
      \end{itemize}
    \end{column}
    \begin{column}{0.5\textwidth}
      \textbf{Atenuação e Doppler}
      \begin{itemize}
        \item alta atenuação em mmWave (drone $\rightarrow$ solo)
        \item efeito Doppler elevado (satélite LEO a 7,5~km/s)
        \item pre-compensação de frequência
      \end{itemize}
      \textbf{Regulatório}
      \begin{itemize}
        \item espectro dedicado vs.\ ISM
        \item licenciamento por país
        \item segurança e privacidade
      \end{itemize}
    \end{column}
  \end{columns}
  \note{O efeito Doppler em LEO (satélite a ~7,6 km/s) pode ultrapassar dezenas de kHz em bandas altas; por isso o 3GPP define a pre-compensação de frequência no enlace de subida.}
\end{frame}

\begin{frame}{Exercício resolvido --- órbita e visibilidade de um satélite LEO}
  \textbf{Problema:} satélite LEO a $h = 550$~km de altitude (faixa do Starlink). Calcule velocidade orbital, período e tempo de visibilidade (elevação mínima $0^\circ$).
  \begin{enumerate}
    \item \textbf{Raio orbital:} $r = R_{\oplus} + h = 6371 + 550 = 6921$~km $= 6{,}921 \times 10^6$~m.
    \item \textbf{Velocidade orbital} ($\mu = 3{,}986 \times 10^{14}$~m$^3$/s$^2$):
      $$v = \sqrt{\dfrac{\mu}{r}} = \sqrt{\dfrac{3{,}986 \times 10^{14}}{6{,}921 \times 10^6}} \approx 7{,}59 \times 10^3\ \text{m/s} \approx 7{,}6\ \text{km/s}.$$
    \item \textbf{Período orbital:}
      $$T = 2\pi\sqrt{\dfrac{r^3}{\mu}},\qquad r^3 = 3{,}315 \times 10^{20}\ \text{m}^3,\quad \dfrac{r^3}{\mu} \approx 831\,705\ \text{s}^2,$$
      $$\sqrt{831\,705} \approx 912\ \text{s}\ \Rightarrow\ T = 2\pi \times 912 \approx 5730\ \text{s} \approx 95{,}5\ \text{min}.$$
    \item \textbf{Tempo de visibilidade} (arco $\eta = \arccos(R_{\oplus}/r) = \arccos(6371/6921) \approx 23{,}0^\circ$):
      $$t_{\text{vis}} = T \times \dfrac{2\eta}{360^\circ} = 5730 \times \dfrac{46}{360} \approx 732\ \text{s} \approx 12{,}2\ \text{min}.$$
  \end{enumerate}
  \begin{block}{Verificação}
    Satélites Starlink (540--570~km) têm período de $\approx 95$--$96$~min, coerente com o resultado.
  \end{block}
\end{frame}

\begin{frame}{Resumo da aula}
  \begin{itemize}
    \item \textbf{FANET}: rede de VANTs com alta mobilidade, topologia 3D, enlaces drone a drone e energia limitada (bateria).
    \item \textbf{NTN}: plataformas aéreas/espaciais; LEO (300--1500~km, 20--40~ms), MEO, GEO ($>$ 250~ms) e HAPS.
    \item \textbf{Constelações LEO}: Starlink ($>$ 7000 sat), OneWeb, Kuiper; ISL a laser e latência competitiva.
    \item \textbf{Integração 5G}: 3GPP Rel.~17 (suporte básico NTN) e Rel.~18 (NTN NR, IoT-NTN).
    \item \textbf{Desafios}: Doppler elevado ($\approx 7{,}6$~km/s em LEO), \emph{handover} entre satélites, atenuação mmWave, energia e regulamentação.
    \item \textbf{Física orbital}: $v = \sqrt{\mu/r}$, $T = 2\pi\sqrt{r^3/\mu}$; LEO a 550~km $\rightarrow T \approx 95{,}5$~min.
  \end{itemize}
\end{frame}

\begin{frame}{Autoavaliação}
  \begin{enumerate}
    \item Qual característica mais distingue uma FANET de uma MANET clássica?
      \begin{itemize}
        \item (a) nós estacionários e topologia fixa
        \item (b) alta mobilidade e topologia 3D em mudança constante
        \item (c) uso exclusivo de enlaces com fio
      \end{itemize}
      \pause
      \textbf{Gabarito:} (b)

    \item Qual categoria de NTN oferece a menor latência fim-a-fim?
      \begin{itemize}
        \item (a) GEO (35~786~km)
        \item (b) MEO (7000--25\,000~km)
        \item (c) LEO (300--1500~km)
      \end{itemize}
      \pause
      \textbf{Gabarito:} (c)

    \item O período orbital de um satélite depende apenas de:
      \begin{itemize}
        \item (a) da inclinação da órbita e da massa do satélite
        \item (b) do raio orbital e da constante gravitacional $\mu$
        \item (c) do tamanho da antena do terminal
      \end{itemize}
      \pause
      \textbf{Gabarito:} (b)

    \item Por que o enlace com satélite LEO sofre Doppler elevado?
      \begin{itemize}
        \item (a) pela alta velocidade relativa do satélite ($\approx 7{,}6$~km/s)
        \item (b) pela baixa altitude em relação ao horizonte
        \item (c) pela banda Ku ser intrinsicamente instável
      \end{itemize}
      \pause
      \textbf{Gabarito:} (a)
  \end{enumerate}
\end{frame}

\begin{frame}{Laboratório sugerido}
  \begin{block}{Visualização de constelações}
    \begin{itemize}
      \item Explorar \emph{Stuff in Space} (stuffin.space) para observar em tempo real a distribuição de Starlink/OneWeb em LEO e a densidade de satélites sobre o Brasil.
      \item Comparar altitude, inclinação e período exibidos com o valor teórico ($T \approx 95$~min a 550~km).
    \end{itemize}
  \end{block}
  \begin{block}{Simulação de FANET no ns-3}
    \begin{itemize}
      \item Configurar um cenário com $N$ UAVs, mobilidade 3D e roteamento ad hoc (ex.: AODV/OLSR).
      \item Medir PDR (packet delivery ratio) e atraso fim-a-fim variando a velocidade dos nós.
      \item Observar o impacto da mobilidade na frequência de trocas de rota e de \emph{handover}.
    \end{itemize}
  \end{block}
  \begin{alertblock}{Atenção}
    Simular em ambiente controlado; não interferir com redes reais sem autorização.
  \end{alertblock}
\end{frame}

\section{Bibliografia}

\begin{frame}{Bibliografia}
  \begin{itemize}
    \item 3GPP TR 38.811, \emph{Study on NR to support NTN}
    \item 3GPP TR 38.821, \emph{Solutions for NR to support NTN}
    \item IEEE Commun. Mag., Special Issue on 5G NTN, 2022
    \item O. S. Oubbati et al., ``A survey on UAV-aided vehicular networks'', \emph{IEEE Commun. Surv. Tut.}, 2020
    \item M. Mozaffari et al., ``A tutorial on UAVs for wireless networks'', \emph{IEEE Commun. Surv. Tut.}, 2019
  \end{itemize}
\end{frame}

\begin{frame}{Leitura recomendada}
  \begin{itemize}
    \item 3GPP TR 38.811 e TR 38.821 --- estudo e soluções para NR-NTN (leitura oficial dos requisitos).
    \item M. Mozaffari et al., ``A tutorial on UAVs for wireless networks'', \emph{IEEE Commun. Surv. Tut.}, 2019 --- taxonomia de usos de UAV como \emph{BS aéreas}, relés e rede de sensores.
    \item O. S. Oubbati et al., ``A survey on UAV-aided vehicular networks'', \emph{IEEE Commun. Surv. Tut.}, 2020 --- integração UAV--VANET.
    \item Especial 5G NTN do \emph{IEEE Commun. Mag.} (2022) --- artigos sobre LEO, MEO e handover satelital.
  \end{itemize}
  \begin{alertblock}{Dica de estudo}
    Releia a Física orbital ($v = \sqrt{\mu/r}$, $T = 2\pi\sqrt{r^3/\mu}$) e confronte com os dados das constelações na \emph{Stuff in Space}.
  \end{alertblock}
\end{frame}

\end{document}